% from ega4-2-fr.tex
\oldpage[IV-2]{182}
On ne peut pas non plus dans cet énoncé remplacer « géométriquement
unibranche » par « unibranche ». Soit en effet $A$ l'anneau local intègre
complet
\[
\mathbb R[[U,V]]/(U^2+V^2);
\]
si $u$, $v$ sont les images de $U$, $V$ dans $A$, l'idéal maximal de $A$
est $Au+Av$. On vérifie aisément que la clôture intégrale de $A$ est
l'anneau $A[t]$, où $t=u/v$ vérifie la relation $t^2=-1$, si bien que
$A[t]$ est isomorphe à l'anneau local $\mathbb C[[U]]$; comme le corps
résiduel de $\mathbb C[[U]]$ est $\mathbb C$, on voit que $A$ est
unibranche mais non géométriquement unibranche. Mais
$A\otimes_{\mathbb R}\mathbb C$ n'est pas un anneau intègre, étant
isomorphe à
\[
\mathbb C[[U,V]]/(U-iV)(U+iV).
\]
\end{remarks}

\section{Relations entre un anneau local noethérien et son complété. Anneaux excellents}
\label{section:IV.7-fr}

Le présent paragraphe est surtout consacré à l'exposé de certaines propriétés
d'anneaux noethériens, généralement stables par passage à une algèbre de type
fini ou par localisation, qui sont vraies pour les anneaux que l'on rencontre
le plus souvent (algèbres de type fini sur $\mathbb Z$, ou sur un corps, ou sur
un anneau local noethérien complet), sans l'être pour tous les anneaux
noethériens. Une première série de propriétés, liées à la théorie de la
dimension, et notamment à la condition des chaînes, fait l'objet des
n\textsuperscript{os} 1 et
2; toutes les propriétés envisagées sont vraies pour les quotients d'anneaux
noethériens réguliers et leurs démonstrations dans ce cas sont pour la plupart
faciles et bien connues [30]; aussi pour ces anneaux les développements des
\S\S~1 et 2 sont sans intérêt.

Il n'en est plus de même des propriétés étudiées dans les
n\textsuperscript{os} 3, 4 et 5 (qui ne dépendent pas des
n\textsuperscript{os} 1 et 2); leurs démonstrations, même dans le cas des
localisés d'algèbres de type fini sur un corps ou sur $\mathbb Z$ (où elles
sont dues pour la plupart à Zariski et à Nagata) sont le plus souvent
délicates; des exemples classiques sont par exemple les équivalences
$A$ normal $\Longleftrightarrow \widehat A$ normal, et $A$ réduit
$\Longleftrightarrow \widehat A$ réduit. Nous développons une méthode
systématique pour formuler et prouver de telles propriétés, à l'aide des
propriétés des \emph{fibres} du morphisme canonique
$\operatorname{Spec}(\widehat A)\to\operatorname{Spec}(A)$, et en utilisant
les notions et résultats exposés au § 6; le succès de cette méthode repose en
dernière analyse sur les excellentes propriétés des \emph{anneaux locaux
complets} et
de ceux qui s'en déduisent par localisation ou passage à une algèbre de type
fini. A cet égard, le théorème de Nagata (6.12.7) disant que, pour un tel
anneau, le lieu singulier dans $\operatorname{Spec}(A)$ est fermé, joue un rôle
crucial; il en est de même (en ce qui concerne les propriétés de permanence)
du critère de régularité (0, 22.3.4), de nature plus technique.

Enfin, dans les n\textsuperscript{os} 6 et 7, nous appliquons les résultats
obtenus à l'étude de
la finitude de la clôture intégrale des anneaux intègres noethériens.

Le lecteur notera que les résultats les plus délicats du § 7 ne serviront
qu'assez peu dans la suite du chapitre IV, et même dans les chapitres
ultérieurs.

\subsection{Équidimensionalité formelle et anneaux formellement caténaires}
\label{subsection:IV.7.1-fr}

